\documentclass[12pt]{article}

\usepackage[a4paper,margin=2.5cm]{geometry}

\usepackage[onehalfspacing]{setspace}

\usepackage{amsmath}
\usepackage{amssymb}
\usepackage{xcolor}
\usepackage{mathtools}
\usepackage{amsthm}
\newtheorem{definition}{Definition}
\newtheorem{example}{Example}

\newenvironment{solution}
{\begin{proof}[Solution]}
{\end{proof}}

\setlength{\parindent}{0pt}

\begin{document}

\begin{titlepage}
    \centering

    \vspace*{\fill}

    {\Huge Linear Transformations I}

    \vspace{1cm}

    {\Large Jin-Ren Ryan Wang\par}

    \vspace{0.5cm}

    {\large \today\par}

    \vspace*{\fill}
\end{titlepage}

\tableofcontents

\newpage

\section{Linear Transformation}

\subsection{Motivation: Matrices as Functions}

Let $F$ be a field and $M_{m\times n}(F)$ be the collection of all $m\times n$ matrices over the field $F$.

Let $A\in M_{m\times n}(F)$ be such a matrix. Every matrix naturally defines a function
\[
f_A:F^n\rightarrow F^m \text{ defined by } f_A(v)=Av.
\]

By using properties of matrix multiplication, we have

\begin{itemize}
    \item $f_A(\mathbf{v_1+v_2})=A(\mathbf{v_1+v_2})=A\mathbf{v_1}+A\mathbf{v_2}=f_A(\mathbf{v_1})+f_A(\mathbf{v_2})$ for all $\mathbf{v_1,v_2}\in F^n$.
    \item $f_A(\lambda \mathbf{v})=A(\lambda \mathbf{v})=\lambda(A\mathbf{v})=\lambda f_A(\mathbf{v})$ for every $\lambda\in F$ and $\mathbf{v}\in F^n$.
\end{itemize}

Our goal is to study functions on \emph{other vector spaces} that \emph{behave like matrices}.\\
We call them \emph{linear transformations}.

\subsection{Definition}

Matrices can be viewed as functions between vector spaces.  
This motivates the following definition.

\textbf{Definition 4.2.1.}
Let $V$ and $W$ be vector spaces over the same field $F$. A function $f:V\to W$ is called a
\textbf{linear transformation}/\textbf{linear map}/\textbf{linear function}
if it satisfies:

\begin{enumerate}
    \item (Additivity) $f(\mathbf{v_1+v_2})=f(\mathbf{v_1})+f(\mathbf{v_2}),\quad\forall\,\mathbf{v_1,v_2}\in V.$

    \item (Homogeneity) $f(\lambda \mathbf{v})=\lambda f(\mathbf{v}),\quad\forall\,\lambda\in F,\mathbf{v}\in V.$
\end{enumerate}

\subsection{Theorem}

\textbf{Theorem 4.3.1.}
Suppose $f:V\rightarrow W$ be a linear transformation. Then $f(\mathbf{0}_V)=\mathbf{0}_W$.

\begin{proof}
Since
\[
\mathbf{0}_V=\mathbf{0}_V+\mathbf{0}_V,
\]
we have
\[
\begin{aligned}
f(\mathbf{0}_V) =f(\mathbf{0}_V+\mathbf{0}_V)=f(\mathbf{0}_V)+f(\mathbf{0}_V) \implies f(\mathbf{0}_V)=\mathbf{0}_W.
\end{aligned}
\]

\end{proof}

\newpage

\textbf{Theorem 4.3.2. (Linear maps are determined by their action on a basis)}

Let $f:V\rightarrow W$ be a linear transformation, where $V$ is finite-dimensional.

Let $\{\mathbf{v}_1,\mathbf{v}_2,\ldots,\mathbf{v}_n\}$ is a basis of $V$. Then $f$ is determined by $f(\mathbf{v}_1),\,f(\mathbf{v}_2),\,\ldots,\,f(\mathbf{v}_n)$.

More precisely, it means

\begin{enumerate}
    \item $\forall \mathbf{v} \in V$, $f(\mathbf{v})$ is a linear combination of $f(\mathbf{v}_1),f(\mathbf{v}_2),\ldots,f(\mathbf{v}_n)$.

    \item Suppose the linear map $g:V\rightarrow W$ satisfies $g(\mathbf{v}_i)=f(\mathbf{v}_i),\quad\forall i=1,\ldots,n$.\\
    Then $g(\mathbf{v})=f(\mathbf{v})$, $\forall \mathbf{v}\in V$.
\end{enumerate}

\begin{proof}

Let $\mathbf{v}\in V$. Since $\{\mathbf{v}_1,\mathbf{v}_2,\ldots,\mathbf{v}_n\}$ is a basis of $V$, we can write

$\mathbf{v}=\alpha_1\mathbf{v}_1+\alpha_2\mathbf{v}_2+\cdots+\alpha_n\mathbf{v}_n$. Then
\[
\begin{aligned}
f(\mathbf{v})=f(\alpha_1\mathbf{v}_1+\cdots+\alpha_n\mathbf{v}_n)=\alpha_1f(\mathbf{v}_1)+\cdots+\alpha_nf(\mathbf{v}_n).
\end{aligned}
\]

In particular, $f(\mathbf{v})$ is determined by $f(\mathbf{v}_1),\ldots,f(\mathbf{v}_n)$. This proves (1).

\medskip

Now suppose another linear transformation $g:V\rightarrow W$ satisfies $g(\mathbf{v}_i)=f(\mathbf{v}_i)$, then
\[
\begin{aligned}
g(\mathbf{v})
&=
g(\alpha_1\mathbf{v}_1+\cdots+\alpha_n\mathbf{v}_n)\\
&=
\alpha_1g(\mathbf{v}_1)+\cdots+\alpha_ng(\mathbf{v}_n)\\
&=
\alpha_1f(\mathbf{v}_1)+\cdots+\alpha_nf(\mathbf{v}_n)\\
&=
f(\alpha_1\mathbf{v}_1+\cdots+\alpha_n\mathbf{v}_n)\\
&=
f(\mathbf{v}).
\end{aligned}
\]

This proves (2)

\end{proof}

\subsection{Example}

\subsubsection{Rotation and Reflection}

Let

\[
\operatorname{Rot}_{\theta}:\mathbb{R}^2\rightarrow\mathbb{R}^2
\]

be defined by

\[
\operatorname{Rot}_{\theta}(\mathbf{v})
=
\begin{pmatrix}
\cos\theta & -\sin\theta\\
\sin\theta & \cos\theta
\end{pmatrix}
\mathbf{v}.
\]

Similarly,

\[
\operatorname{Ref}_{\theta}:\mathbb{R}^2\rightarrow\mathbb{R}^2
\]

is defined by

\[
\operatorname{Ref}_{\theta}(\mathbf{v})
=
\begin{pmatrix}
\cos(2\theta) & \sin(2\theta)\\
\sin(2\theta) & -\cos(2\theta)
\end{pmatrix}
\mathbf{v},
\]

which represents reflection along the line

\[
y=\tan(\theta)x.
\]

Since both functions are of the form

\[
f(\mathbf{v})=A\mathbf{v},
\]

they satisfy the two properties of linear transformations.

Hence,

\[
\operatorname{Rot}_{\theta},\;
\operatorname{Ref}_{\theta}
:\mathbb{R}^2\rightarrow\mathbb{R}^2
\]

are both linear transformations.

\subsubsection{Differentiation}

Define $D:\mathbb{R}[x]_{\le n} \longrightarrow \mathbb{R}[x]_{\le n-1}$
by
\[
D(f(x))=f'(x).
\]

To verify that $D$ is linear, check the two defining properties.

\begin{enumerate}
    \item For any $f_1(x),f_2(x)\in\mathbb{R}[x]_{\le n}$,
    \[
    \begin{aligned}
    D(\textcolor{red}{f_1(x)+f_2(x)})
    &=\textcolor{red}{(f_1(x)+f_2(x))'} \\
    &=\textcolor{red}{f_1'(x)+f_2'(x)} \\
    &=D(\textcolor{red}{f_1(x)})+D(\textcolor{red}{f_2(x)}).
    \end{aligned}
    \]

    \item For every $\lambda\in\mathbb{R}$,
    \[
    \begin{aligned}
    D(\lambda \textcolor{red}{f(x)})
    &=(\lambda \textcolor{red}{f(x)})'\\
    &=\lambda \textcolor{red}{f'(x)}\\
    &=\lambda D(\textcolor{red}{f(x)}).
    \end{aligned}
    \]
\end{enumerate}

Therefore, $D$ "behaves like matrices" $\implies$ D is a linear transformation.

\newpage

\subsubsection{Trace}

Recall the trace map
\[
\operatorname{tr}:M_n(\mathbb{C})\rightarrow\mathbb{C},
\]

defined by
\[
\operatorname{tr}(A)=a_{11}+a_{22}+\cdots+a_{nn},
\]

where $A=(a_{ij}).$

\medskip

To verify linearity,
\begin{enumerate}
    \item For any $A,B\in M_n(\mathbb{C})$,
    \[
    \operatorname{tr}(\textcolor{red}{A+B}) = \operatorname{tr}(\textcolor{red}{A}) + \operatorname{tr}(\textcolor{red}{B}).
    \]

    \item For every $\lambda\in\mathbb{C}$,
    \[
    \operatorname{tr}(\lambda \textcolor{red}{A}) = \lambda\operatorname{tr}(\textcolor{red}{A}).
    \]
\end{enumerate}

Hence,
\[
\operatorname{tr}:M_n(\mathbb{C})\rightarrow\mathbb{C}
\]

is a linear transformation.

\subsubsection{Composition Operator}

Fix a continuous function $g\in C(\mathbb{R}).$. Then $S_g:C(\mathbb{R})\rightarrow C(\mathbb{R})$ is defined by
\[
S_g(f(x)) = f(g(x)) \text{ is a linear transformation. }
\]
In particular, if $g(x)=x-\alpha$ (with fixed $\alpha \in \mathbb{R}$), then $S_g$ refers to the function
that shifts the the graph of $f(x)$ by $\alpha$ units to the right.

\begin{solution}

To verify linearity,

\begin{enumerate}
    \item For any $f_1,f_2\in C(\mathbb{R})$,
    \[
    \begin{aligned}
    S_g(\textcolor{red}{f_1+f_2})
    &= \textcolor{red}{f_1(g(x))}+\textcolor{red}{f_2(g(x))}\\
    &= S_g(\textcolor{red}{f_1})+S_g(\textcolor{red}{f_2}).
    \end{aligned}
    \]

    \item For every $\lambda\in\mathbb{R}$,
    \[
    \begin{aligned}
    S_g(\lambda \textcolor{red}{f})
    &= \lambda \textcolor{red}{f(g(x))}\\
    &= \lambda S_g(\textcolor{red}{f}).
    \end{aligned}
    \]
\end{enumerate}

Hence, $S_g:C(\mathbb{R})\rightarrow C(\mathbb{R})$ is a linear transformation.

\end{solution}

\section{Kernel and Image}

\subsection{Definition 4.4.1.}
Let $f:V\rightarrow W$ be a linear transformation. Define
\begin{enumerate}
    \item the \textbf{kernel} (or \textbf{nullspace}) of $f$ by
    \[
    \operatorname{ker}(f)
    =
    \left\{
    \mathbf{v}\in V
    \mid
    f(\mathbf{v})=\mathbf{0}_W
    \right\}.
    \]

    i.e. the kernel consists of all vectors in the domain that become the zero vector under $f$.

    \item the \textbf{image} (or \textbf{range}) of $f$ by
    \[
    \operatorname{im}(f)
    =
    \left\{
    \mathbf{w}\in W
    \mid
    \mathbf{w}=f(\mathbf{v})
    \text{ for some }
    \mathbf{v}\in V
    \right\}.
    \]

    i.e. the image/range consists of all the possible outputs of the function $f$.
\end{enumerate}

\subsection{Theorem}

\textbf{Theorem 4.4.1. (Properties of kernel)}

Let $f:V\rightarrow W$ be a linear transformation.

\begin{enumerate}
    \item $\operatorname{ker}(f)$ is a subspace of $V$.

    \item $f$ is injective if and only if $\operatorname{ker}(f)=\{\mathbf{0}_V\}$.
\end{enumerate}

\begin{proof}

(1) To show that $\operatorname{ker}(f)$ is a subspace of $V$, we verify the three subspace conditions.

\begin{itemize}
    \item \fbox{Nonempty.}

    Since $f$ is linear, $f(\mathbf{0}_V)=\mathbf{0}_W$ $\implies \mathbf{0}_V\in\operatorname{ker}(f).$

    \item \fbox{Closed under addition.}

    Let $\mathbf{v}_1,\mathbf{v}_2\in\operatorname{ker}(f)$ $\implies f(\mathbf{v}_1)=\mathbf{0}_W,\quad f(\mathbf{v}_2)=\mathbf{0}_W.$
    
    Check:
    \[
    \begin{aligned}
    f(\mathbf{v}_1+\mathbf{v}_2)
    &=
    f(\mathbf{v}_1)+f(\mathbf{v}_2)\\
    &=
    \mathbf{0}_W+\mathbf{0}_W\\
    &=
    \mathbf{0}_W.
    \end{aligned}
    \]
    $\implies \mathbf{v}_1+\mathbf{v}_2\in\operatorname{ker}(f)$.
    \newpage

    \item \fbox{Closed under scalar multiplication.}

    Let
    $\lambda\in F$ and $\mathbf{v}\in\operatorname{ker}(f)$.

    Then
    \[
    \begin{aligned}
    f(\lambda\mathbf{v})
    &=
    \lambda f(\mathbf{v})\\
    &=
    \lambda\mathbf{0}_W\\
    &=
    \mathbf{0}_W.
    \end{aligned}
    \]
    $\implies \lambda\mathbf{v}\in\operatorname{ker}(f)$.

\end{itemize}

Therefore, $\operatorname{ker}(f)$ is a subspace of $V$.

\medskip

(2) We show both directions.

$(\Rightarrow)$

Suppose $f$ is injective. If $\mathbf{v}\in\operatorname{ker}(f)$, then

\[
f(\mathbf{v})=\mathbf{0}_W=f(\mathbf{0}_V).
\]

Since $f$ is injective, $\mathbf{v}=\mathbf{0}_V$ $\implies \operatorname{ker}(f) = \{\mathbf{0}_V\}$.

$(\Leftarrow)$

Suppose $\operatorname{ker}(f)=\{\mathbf{0}_V\}$. If $f(\mathbf{v}_1)=f(\mathbf{v}_2)$, then

\[
\begin{aligned}
f(\mathbf{v}_1-\mathbf{v}_2)= f(\mathbf{v}_1)-f(\mathbf{v}_2)=\mathbf{0}_W.\\
\implies \mathbf{v}_1-\mathbf{v}_2 \in \ker(f)   
\end{aligned}
\]

Therefore, $\mathbf{v}_1-\mathbf{v}_2=\mathbf{0}_V$, which implies $\mathbf{v}_1=\mathbf{v}_2$.

Thus $f$ is injective.

\end{proof}

\textbf{Theorem 4.4.2. (Properties of image)}

Let $f:V\rightarrow W$ be a linear transformation.

\begin{enumerate}
    \item $\operatorname{im}(f)$ is a subspace of $W$.

    \item $f$ is surjective if and only if $\operatorname{im}(f)=W$. \textcolor{red}{($\leftarrow$ def of surjectivity (nothing to prove))}
\end{enumerate}

\begin{proof}

(1) We verify the three subspace conditions.

\begin{itemize}
    \item \fbox{Nonempty.}

    Since $f(\mathbf{0}_V)=\mathbf{0}_W$, we have $\mathbf{0}_W\in\operatorname{im}(f)$.

    \item \fbox{Closed under addition.}

    Let $\mathbf{w}_1,\mathbf{w}_2 \in \operatorname{im}(f)$.

    Then there $\exists$ $\mathbf{v}_1,\mathbf{v}_2\in V$ s.t. $\mathbf{w}_1=f(\mathbf{v}_1),\quad\mathbf{w}_2=f(\mathbf{v}_2)$.

    Hence,
    \[
    \begin{aligned}
    \mathbf{w}_1+\mathbf{w}_2
    &=
    f(\mathbf{v}_1)+f(\mathbf{v}_2)\\
    &=
    f(\mathbf{v}_1+\mathbf{v}_2).
    \end{aligned}
    \]

    Therefore, $\mathbf{w}_1+\mathbf{w}_2\in\operatorname{im}(f)$.

    \item \fbox{Closed under scalar multiplication.}

    Let $\mathbf{w}=f(\mathbf{v})$. Then $\lambda\mathbf{w}=\lambda f(\mathbf{v})=f(\lambda\mathbf{v})$, so $\lambda\mathbf{w}\in\operatorname{im}(f)$.
\end{itemize}

Therefore, $\operatorname{im}(f)$ is a subspace of $W$.

\medskip

(2)

By definition,

\[
\begin{aligned}
f \text{ is surjective}
&\iff
\forall\,\mathbf{w}\in W,\;
\exists\,\mathbf{v}\in V
\text{ s.t. } f(\mathbf{v})=\mathbf{w}\\
&\iff
W\subseteq\operatorname{im}(f).   
\end{aligned}
\]

Since $\operatorname{im}(f)\subseteq W$ always holds, we conclude
\[
f
\text{ is surjective}
\iff
\operatorname{im}(f)=W.
\]

\end{proof}

\subsection{Example}

\subsubsection{Differentiation}

Consider $D:\mathbb{R}[x]_{\le3}\longrightarrow \mathbb{R}[x]_{\le2}$ defined by
\[
D(f(x))=f'(x).
\]

\begin{solution}
First, compute the kernel.

\[
\begin{aligned}
\operatorname{ker}(D)
&=
\{\textcolor{red}{f\in\mathbb{R}[x]_{\le3}}:D(\textcolor{red}{f})=0\}\\
&=
\{f\in\mathbb{R}[x]_{\le3}:\textcolor{red}{f'(x)}=0\}\\
&=
\{\text{constant polynomials}\}\\
&=\mathbb{R}.
\end{aligned}
\]

As $\operatorname{ker}(D)\neq\{\mathbf{0}\}$, we deduce that $D$ is not injective.

Next, compute the image. Let
\[
p(x)=ax^2+bx+c\in \mathbb{R}[x]_{\le2}.
\]

Then
\[
p(x)=D\!
\left(\textcolor{red}{\frac{a}{3}x^3+\frac{b}{2}x^2+cx}\right).
\]

As $\operatorname{im}(D)=\mathbb{R}[x]_{\le2}$ (whole domain), $D$ is surjective.

\end{solution}

\subsubsection{Trace}

Consider $\operatorname{tr}:M_2(\mathbb{R})\longrightarrow\mathbb{R}$ defined by
$
\operatorname{tr}
\begin{pmatrix}
a&b\\
c&d
\end{pmatrix}
= a+d.
$

\begin{solution}
First, compute the kernel.
\[
\begin{aligned}
\operatorname{ker}(\operatorname{tr})
&=
\left\{
A\in M_2(\mathbb{R}):\operatorname{tr}(A)=0
\right\}\\
&=
\left\{
\begin{pmatrix}
a&b\\
c&\textcolor{red}{-a}
\end{pmatrix}
:
a,b,c\in\mathbb{R}
\right\}.
\end{aligned}
\]

Equivalently,
\[
\operatorname{ker}(\operatorname{tr})=\operatorname{span}
\left\{
\begin{pmatrix}
1&0\\
0&-1
\end{pmatrix},
\begin{pmatrix}
0&1\\
0&0
\end{pmatrix},
\begin{pmatrix}
0&0\\
1&0
\end{pmatrix}
\right\}.
\]

Hence, $\operatorname{ker}(\operatorname{tr})\neq\{\mathbf{0}\}$, so $\operatorname{tr}$ is not injective.

Next, compute the image. For every $r\in\mathbb{R}$,
\[
r=\operatorname{tr}
\begin{pmatrix}
r&0\\
0&0
\end{pmatrix}.
\]

Thus, $\operatorname{im}(\operatorname{tr})=\mathbb{R}$ $\implies\operatorname{"tr"}$ is surjective.

\end{solution}

\section{Rank-Nullity Theorem}

\subsection{Definition}

\textbf{Definition 4.5.1.}
Let $f: V \rightarrow W$ be a linear map. Define
\begin{itemize}
    \item[1.] the \emph{nullity} of $f$ by $\operatorname{nullity}(f)$ = $\operatorname{dim}(\operatorname{ker}(f))$.
    \item[2.] the \emph{rank} of $f$ by $\operatorname{rank}(f)$ = $\operatorname{dim}(\operatorname{ker}(f))$
\end{itemize}

\newpage

\subsection{Theorem}

\textbf{Theorem 4.5.1. (Rank-Nullity Theorem)}

Let $V$ be a finite-dimensional vector space, and let $f:V\rightarrow W$ be a linear transformation.

Then we have
\[
\operatorname{nullity}(f)+\operatorname{rank}(f)=\operatorname{dim}(V).
\]


\begin{proof}
\fbox{Idea: $\operatorname{ker}(f)\subseteq V$}

Let $S$ = $\{\mathbf{v}_1,\mathbf{v}_2,\ldots,\mathbf{v}_n\}$ be a basis of $\operatorname{ker}(f)\quad$
(i.e. $\operatorname{nullity}(f)$ = $\operatorname{dim}(\operatorname{ker}(f))$)

By the Basis Extension Theorem, we may extend $S$to a basis of $V$:
\[
S'=\{\mathbf{v}_1,\ldots,\mathbf{v}_n,\textcolor{red}{\mathbf{s}_1,\ldots,\mathbf{s}_k}\}\text{ : a basis of }V.\text{ (i.e. $\operatorname{dim}(V)=n+k$)} 
\]

We claim that $S''=\{\textcolor{red}{f(\mathbf{s}_1),\ldots,f(\mathbf{s}_k)}\}$ forms a basis of $\operatorname{im}(f)$.

\begin{proof}[Proof of Claim]

We first show that it spans $\operatorname{im}(f)$.

Let $\mathbf{w}\in\operatorname{im}(f)$ $\implies \mathbf{w}=f(\textcolor{red}{\mathbf{v}}) \text{ for \textcolor{red}{some $\mathbf{v} \in V$} } $

As $S'$ is a basis of $V$, we can write
\[
\textcolor{red}{\mathbf{v}}=\sum_{i=1}^{n}\alpha_i\mathbf{v}_i+\sum_{i=1}^{k}\beta_i\textcolor{red}{\mathbf{s}_i}.
\]

Then,
\begin{align*}
\mathbf{w}
&=
f(\textcolor{red}{\mathbf{v}}) \\
&=
f\!\left(
\sum_{i=1}^{n}\alpha_i\mathbf{v}_i+\sum_{i=1}^{k}\beta_i\textcolor{red}{\mathbf{s}_i}
\right) \\
&=
\sum_{i=1}^{n}\alpha_i\textcolor{red}{f(\mathbf{v}_i)}+\sum_{i=1}^{k}\beta_i\textcolor{red}{f(\mathbf{s}_i)} \\
&= 
\sum_{i=1}^{k}\beta_i\textcolor{red}{f(\mathbf{s}_i)} \implies S''\text{ spans }\operatorname{im}(f)
\end{align*}

Next, we prove linear independence. Suppose
\[
\sum_{i=1}^{k}r_i\textcolor{red}{f(\mathbf{s}_i)=\mathbf{0}}.
\]

Then
\[
\textcolor{red}{f}\!\left(
\sum_{i=1}^{k}r_i\mathbf{s}_i
\right)
= \mathbf{0},
\]

so
\[
\sum_{i=1}^{k}r_i\mathbf{s}_i\in\operatorname{ker}(f).
\]

As $S$ is a basis of $\operatorname{ker(f)}$,
\[
\sum_{i=1}^{k}r_i\mathbf{s}_i=\sum_{i=1}^{n}f_i\mathbf{v}_i \text{ for some scalars $f_i$. }
\]

$\implies\sum_{i=1}^{k}r_i\mathbf{s}_i+\sum_{i=1}^{n}(-f_i)\mathbf{v}_i=\mathbf{0}$

\medskip

As $S'$ is a basis of $V$, it is linearly independent, this forces $r_i=0, f_i = 0$, $\forall i$.

Thus, $\{f(\mathbf{s}_1),\ldots,f(\mathbf{s}_k)\}$ is linearly independent.Hence it forms a basis of $\operatorname{im}(f)$.

\end{proof}

Consequently,
\[
\operatorname{rank}(f)=\operatorname{dim}(\operatorname{im}(f))=\textcolor{red}{k}=(n+k)-n=\operatorname{dim(V)}-\operatorname{nullity}(f),
\]

\end{proof}

\subsection{Example}

\subsubsection{Differentiation}

Recall the linear transformation $D:\mathbb{R}[x]_{\le3}\longrightarrow \mathbb{R}[x]_{\le2}$ defined by
\[
D(f(x))=f'(x).
\]

We have $\operatorname{ker}(D)=\mathbb{R}\text{ and }\operatorname{im}(D)=\mathbb{R}[x]_{\le2}$.
\begin{solution}
clearly,
\[
\operatorname{nullity}(D)= \operatorname{dim}_{\textcolor{red}{\mathbb{R}}}(\operatorname{ker}(D))=\operatorname{dim}_{\textcolor{red}{\mathbb{R}}}(\mathbb{R})=1
\]   
and
\[
\operatorname{rank}(D)=\operatorname{dim}_{\textcolor{red}{\mathbb{R}}}(\operatorname{im}(D))=\operatorname{dim}_{\textcolor{red}{\mathbb{R}}}(\mathbb{R}[x]_{\le2})=3
\] 
\end{solution}

\newpage

\subsubsection{Trace}

Recall that for tr: $\operatorname{tr}:M_2(\mathbb{R})\longrightarrow \mathbb{R}$ defined by
\[
\operatorname{tr}
\begin{pmatrix}
a&b\\
c&d
\end{pmatrix}
= 
a+d,
\]

we have
$
\operatorname{ker}(\operatorname{tr})=\operatorname{span}
\left\{
\begin{pmatrix}
1&0\\
0&-1
\end{pmatrix},
\begin{pmatrix}
0&1\\
0&0
\end{pmatrix},
\begin{pmatrix}
0&0\\
1&0
\end{pmatrix}
\right\}
$
and $\operatorname{im}(\operatorname{tr})=\mathbb{R}$.

\begin{solution}
we have
\[
\operatorname{nullity}(\operatorname{tr})=\operatorname{dim}(\operatorname{ker}(\operatorname{tr}))=3,
\]

and
\[
\operatorname{rank}(\operatorname{tr})=\operatorname{dim}(\operatorname{im}(\operatorname{tr}))=\operatorname{dim}(\mathbb{R})=1.
\]

Notice that
\[
3+1=\textcolor{red}{4=\operatorname{dim}(\underbrace{M_2(\mathbb{R})}_{\text{domain}})},
\]

which equals the dimension of the domain.    
\end{solution}


\subsubsection{Identity Map}
Let $V$ be an $n$-dimensional vector space, and consider the identity map $\operatorname{id}_V:V\rightarrow V$ defined by

\[
\operatorname{id}_V(\mathbf{v})=\mathbf{v},\quad\forall\,\mathbf{v}\in V.
\]

\begin{solution}

Note that:
\[
\operatorname{ker}(\operatorname{id}_V)=\{\mathbf{0}_V\},\quad\operatorname{im}(\operatorname{id}_V)=V,
\]

we obtain
\[
\operatorname{nullity}(\operatorname{id}_V)=\operatorname{dim}(\operatorname{ker}(\operatorname{id}_V))=\operatorname{dim}(\{\mathbf{0}_V\})=0,
\]

and
\[
\operatorname{rank}(\operatorname{id}_V)=\operatorname{dim}(\operatorname{im}(\operatorname{id}_V))=\operatorname{dim}(V)=\textcolor{red}{\text{dimension of the domain of }f}=n.
\]

Hence,
\[
\operatorname{rank}(\operatorname{id}_V)+\operatorname{nullity}(\operatorname{id}_V)=n=\operatorname{dim}(V).
\]

\end{solution}

\newpage

\section{Rounding up}

\subsection{Remark}

\begin{enumerate}

\item (Nullity and Rank of a Matrix)

A matrix $A\in M_{m\times n}(F)$ naturally induces the linear transformation $f_A:F^n\rightarrow F^m$ defined by
$f_A(\mathbf{v})=A\mathbf{v}$. From now on, we identify the matrix $A$ with its associated linear transformation $f_A$.
and denote $f_A$ simply by $A$. That is: 
\[
\operatorname{ker}(A)=\operatorname{ker}(f_A) \text{ and } \operatorname{im}(A)=\operatorname{im}(f_A),
\]

\item (Injectivity and Surjectivity)

For a linear transformation $f:V\rightarrow W$, we have
\[
\operatorname{nullity}(f)=0
\iff
\operatorname{ker(f)}=\{\mathbf{0}_V\}
\iff
f\text{ is injective}.
\]

Dually,
\[
\operatorname{rank}(f)=\operatorname{dim}(W)
\iff
\operatorname{im}(f)=W
\iff
f\text{ is surjective}.
\]

\end{enumerate}


\end{document}